\section{Discussion}\label{sec:discussion}

\subsection{What the results establish, and what they do not}
\emph{Established under stated assumptions:} chunked and parallel processing introduce no detection difference relative to the whole-document pipeline, for detectors that respect their declared bounds (Section~\ref{sec:formal}). \emph{Measured:} leaks caused by naive and overlapping chunking on every corpus, including real system logs and court judgments, and none with SafeStream; flat working memory across a 256-fold range of input sizes; near-linear scaling to four cores; the robust detection margins of Section~\ref{sec:results}; and, on real logs, complete IP-address redaction, user names redacted by the opt-in tier on held-out data, and no false alarms from the pattern tiers where Presidio's statistical configurations flagged most lines; on real court judgments, lower personal-data recall than Presidio for the evaluated configurations, and, for a free-text configuration designed on development data, higher character F1 and far higher direct-identifier protection than Presidio on 1,014 held-out judgments, and higher character F1 on a second genre, Wikipedia biographies; on 14,626 real credentials, more fully redacted than gitleaks, with more false alarms on lines people judged harmless. \emph{Not established:} effectiveness on credentials outside CredData's candidate set and file types; document-level protection adequate for public release; robustness to adversarial obfuscation.

\subsection{When is the approach useful despite a low no-leak rate?}
Bounded memory and equivalence do not make redaction complete; they remove one specific failure mode, the leak introduced by chunking itself, and make large inputs processable at all. The value of that depends on the detection coverage for the values that matter in a given deployment. The evidence supports three uses and argues against a fourth.
\begin{enumerate}
\item \emph{Defence in depth for structured secrets in machine-generated text.} Where the sensitive values are credentials and structured identifiers with recognisable shapes (logs, configuration dumps, continuous-integration output), the pattern tiers are where coverage is highest in our data, and the guarantee ensures that splitting a multi-gigabyte log does not make coverage worse. Even here more than half of the test secrets were not fully redacted, so the tool reduces exposure rather than eliminating it.
\item \emph{Pre-filtering before human review or controlled sharing.} When output is reviewed, or shared only with parties already bound by agreements, a high-precision filter that reliably processes large inputs reduces the volume of exposed values without being the last line of defence.
\item \emph{Inputs that cannot otherwise be processed.} Whole-document redaction needed about 28 times the input size in memory in our measurements, which rules it out for multi-gigabyte files on ordinary hosts.
\item \emph{Not supported: sole control for public release of free text.} With no-leak rates below 0.18 on RedactionBench, and a direct identifier still exposed somewhere in about one held-out TAB document--annotator pair in five even with the free-text configuration, the results do not justify releasing redacted free-text documents publicly on the strength of this system alone; with that configuration the system is a strong first pass for prose that a person then reviews.
\end{enumerate}
A deployment should set its own acceptance criteria before relying on the system, for example: per-category recall measured on a representative, authorised sample of the target artefacts; zero tolerated leaks for specified credential types on that sample; a review step for document classes with free text; and periodic re-measurement as formats change. None of these was measured here.

\subsection{Threats to validity}
\emph{External validity.} presidio-research is synthetic, and its result does not carry over to TAB (Section~\ref{sec:tab}); the secrets positives are rule-author examples; RedactionBench is contaminated; all corpora are English. The real-log labels come from templates checked by one author rather than double annotation (a blind second-annotation sheet and agreement script are released, \texttt{benchmarks/\allowbreak loghub\_\allowbreak agreement.py}), and the user-name tier's OpenSSH result is exposed. CredData's labels cover only candidates that other tools proposed, gitleaks and TruffleHog among them; its values are obfuscated; and each document is one line range rather than a whole file. The free-text configuration was designed on TAB and tested on one other genre, where its direct-identifier advantage did not hold against Presidio lg; TAB's test split motivated it and is reported as exposed. \emph{Construct validity.} The secrets gold spans are Nosey Parker's capture groups, so credentials adjacent to a labelled secret count as false positives; RedactionBench's contextual labels have low human agreement \cite{brynjolfsson2026redactionbench}. \emph{Internal validity.} The secrets test half was scored before and after one fix (Section~\ref{sec:fp}). \emph{Statistical validity.} Eleven RedactionBench domains are few clusters for a bootstrap; the ablation intervals are document-level only; no multiplicity correction was applied. \emph{Comparison fairness.} With their filters disabled, gitleaks closed most of the character-F1 gap on the secrets corpus (Section~\ref{sec:sens}), so that margin is not claimed. \emph{Environment.} All runs used one virtual machine.

\subsection{Evaluation plan for real-world evidence}\label{sec:plan}
TAB now provides real free text with independent human gold; CredData provides real credentials; the most valuable next evidence is free text from other domains and real credential-bearing text found by a process independent of existing scanners with clear provenance, such as public continuous-integration logs whose licences permit research use or logs from an organisation that authorises the study, annotated by two people with agreement reported. Safeguards: never test discovered credentials against live services, keep raw text encrypted and limited to annotators, report live secrets to their owners, and publish only annotations, hashes and aggregate results. The repository contains a collection plan, annotation guidelines and a scoring pipeline for such a corpus (\texttt{benchmarks/realworld/}). Transformer-based detectors would additionally need an empirical streamed-versus-whole-document check, because the theorems do not cover a model whose output depends on its whole input window.
